\section{Câu 4.5}
Section 4.2.2 describes a method for generating association rules from frequent itemsets. Propose a more efficient method. Explain why it is more efficient than the one proposed there. (Hint: consider incorporating the properties of Exercises 4.3(b), (c) into your design.)

\paragraph{Lời giải}

Sử dụng tính chất pruning của confidence:

\begin{itemize}
    \item Nếu một luật $s \Rightarrow (l-s)$ không đạt \textit{minconf}, thì mọi luật tạo từ các tập con của $s$ cũng sẽ không đạt \textit{minconf}.
    \item Vì vậy, ta không cần sinh toàn bộ các tập con của $l$.
\end{itemize}

\begin{algorithm}[H]
    \caption{Cải tiến thuật toán sinh luật kết hợp}
    \begin{algorithmic}[1]
        \Require Tập phổ biến $l$, minimum confidence $minconf$
        \ForAll{nonempty subsets $s \subset l$}
        \State $conf \gets \dfrac{support(l)}{support(s)}$

        \If{$conf \ge minconf$}
        \State Output rule:
        \[
            s \Rightarrow (l - s)
        \]
        \State Continue exploring subsets of $s$
        \Else
        \State Prune all subsets of $s$
        \EndIf
        \EndFor
    \end{algorithmic}
    \label{alg:strong_rules}
\end{algorithm}
Phương pháp này hiệu quả hơn vì nó loại bỏ sớm các tập ứng viên không thể tạo ra luật mạnh. Nếu một luật không đạt \textit{minconf}, mọi mở rộng của nó cũng không cần xét tiếp. Do đó, số lượng tập con cần sinh và kiểm tra giảm đáng kể so với cách brute-force sinh toàn bộ $2^{|l|}-2$ tập con.
